\begin{problem}[Legacy problem: generic point of a closed irreducible subset]
Let $X$ be a scheme and let $Z\subseteq X$ be a nonempty closed irreducible subset.
Prove that there exists a unique point $\eta\in Z$ such that
\[
\overline{\{\eta\}}=Z.
\]
In particular, an integral scheme has a unique generic point.
\end{problem}

\paragraph{Why this problem matters.}
This is the fundamental sobriety theorem for schemes.  Generic points are not merely
available in affine spectra: they glue correctly and exist for every irreducible closed
subset of every scheme.

\paragraph{Useful ideas before starting.}
Choose an affine open $U=\Spec A$ meeting $Z$.  The subset $Z\cap U$ is nonempty,
irreducible, and closed in $U$, so it has the form $V(\mathfrak p)$ for a prime
$\mathfrak p$.  Then compare the closure of $\mathfrak p$ in $U$ with its closure in $X$.

\paragraph{Guided hints.}
Let $W=\overline{\{\mathfrak p\}}^{\,X}$.  Since $Z$ is closed and contains
$\mathfrak p$, one has $W\subseteq Z$.  Show
\[
W\cap U=Z\cap U.
\]
If $W\neq Z$, use irreducibility of $Z$ to obtain two disjoint nonempty open subsets of
$Z$.

\begin{solution}
Choose any $x\in Z$ and an affine open neighborhood
\[
U=\Spec A\subseteq X
\]
of $x$.  Then $Z\cap U$ is nonempty, open in $Z$, closed in $U$, and irreducible.
Therefore
\[
Z\cap U=V_U(\mathfrak p)
\]
for a unique prime ideal $\mathfrak p\subseteq A$.  Let $\eta\in U$ be the corresponding
point.  Inside $U$,
\[
\overline{\{\eta\}}^{\,U}=V_U(\mathfrak p)=Z\cap U.
\]

Let
\[
W=\overline{\{\eta\}}^{\,X}.
\]
Because $\eta\in Z$ and $Z$ is closed in $X$, we have $W\subseteq Z$.  Since $U$ is open
and contains $\eta$,
\[
W\cap U
=
\overline{\{\eta\}}^{\,U}
=
Z\cap U.
\]
Suppose $W\neq Z$.  Then $Z\setminus W$ is a nonempty open subset of the irreducible space
$Z$.  The set $Z\cap U$ is also a nonempty open subset of $Z$, but
\[
(Z\setminus W)\cap(Z\cap U)=\varnothing
\]
because $Z\cap U=W\cap U\subseteq W$.  This contradicts irreducibility.  Hence $W=Z$,
so $\eta$ is a generic point of $Z$.

For uniqueness, suppose $\eta'$ is another point with
\[
\overline{\{\eta'\}}=Z=\overline{\{\eta\}}.
\]
Schemes are $T_0$: if two points are distinct, either an affine neighborhood contains one
but not the other, or both lie in an affine spectrum where the $T_0$ property was proved
in this chapter.  In a $T_0$ space distinct points have distinct closures.  Hence
$\eta=\eta'$.

If $X$ is integral, then $X$ itself is irreducible and closed in itself, so it has a unique
generic point.
\end{solution}

\paragraph{What happened mathematically.}
Affine generic points are compatible across overlaps.  This makes a scheme a sober
topological space: irreducible closed geometry can always be represented by a distinguished
point.

\paragraph{Extensions.}
Show that every nonempty open subset of an integral scheme contains its generic point.

\paragraph{Provenance.}
Legacy \emph{Theory of Algebraic Geometry 8}, Problem 1, preserved in full.  The
generic-point portion of legacy \emph{Theory of Algebraic Geometry 6}, Problem 7, is merged
here because it is the special case $Z=X$ for an integral scheme.  That legacy problem also
asks for the stalk at the generic point and the function field; that distinct mathematical
part is explicitly moved to VI/27, \emph{Integral Schemes and Function Fields}, so it is
not duplicated here.
